\documentclass{article}

\usepackage{amsmath,amssymb}
\usepackage{xurl}
\usepackage[hidelinks]{hyperref}

% Shared commands for notes in docs/paper-gaps/.

\DeclareUnicodeCharacter{2080}{\ifmmode{}_{0}\else\textsubscript{0}\fi}
\DeclareUnicodeCharacter{2081}{\ifmmode{}_{1}\else\textsubscript{1}\fi}
\DeclareUnicodeCharacter{2082}{\ifmmode{}_{2}\else\textsubscript{2}\fi}
\DeclareUnicodeCharacter{2099}{\ifmmode{}_{n}\else\textsubscript{n}\fi}

\newcommand{\Lean}{\textsc{Lean}}
\ExplSyntaxOn
\NewDocumentCommand{\leanid}{m}{
  \ifmmode
    \textsf{\detokenize{#1}}
  \else
    \group_begin:
    \urlstyle{sf}
    \tl_set:Nn \l_tmpa_tl {#1}
    \tl_replace_all:Nnn \l_tmpa_tl {\_} {_}
    \exp_args:NV \path \l_tmpa_tl
    \group_end:
  \fi
}
\ExplSyntaxOff
\providecommand{\tr}{\operatorname{tr}}
\newcommand{\tnleanIssueBase}{https://github.com/LionSR/TNLean/issues/}
\newcommand{\tnleanPrBase}{https://github.com/LionSR/TNLean/pull/}
\newcommand{\tnleanDocBase}{https://sirui-lu.com/TNLean/docs/}
\newcommand{\ghissue}[1]{\href{\tnleanIssueBase#1}{GitHub issue \##1}}
\newcommand{\ghpr}[1]{\href{\tnleanPrBase#1}{PR \##1}}
\newcommand{\doclink}[1]{\href{\tnleanDocBase#1.html}{\path{#1}}}

% Dirac notation and the identity operator, as in blueprint/src/macros/common.tex.
% \providecommand keeps note-local definitions authoritative.
\usepackage{dsfont}
\providecommand{\ket}[1]{|#1\rangle}
\providecommand{\bra}[1]{\langle #1|}
\providecommand{\braket}[2]{\langle #1 | #2 \rangle}
\providecommand{\ketbra}[2]{|#1\rangle\!\langle #2|}
\providecommand{\Id}{\mathds{1}}
\providecommand{\one}{\mathds{1}}
\providecommand{\C}{\mathbb{C}}
\providecommand{\MN}[1]{M_{#1}(\C)}
\providecommand{\MD}{M_D(\C)}

% External referencing for paper sources.  The build helper writes sanitized
% auxiliary files under build/paper-aux/ and excludes duplicated source labels.
\usepackage{xr}

% Import a generated paper auxiliary file with a caller-supplied label prefix.
% The candidate paths support builds from docs/paper-gaps/, the repository root,
% and build/paper-aux/ itself.
\newcommand{\importpaperaux}[2]{%
  \IfFileExists{../../build/paper-aux/#2.aux}{%
    \externaldocument[#1]{../../build/paper-aux/#2}%
  }{%
    \IfFileExists{build/paper-aux/#2.aux}{%
      \externaldocument[#1]{build/paper-aux/#2}%
    }{%
      \IfFileExists{#2.aux}{%
        \externaldocument[#1]{#2}%
      }{}%
    }%
  }%
}

% General external reference macros
\newcommand{\paperref}[2]{\ref{#1-#2}}
\newcommand{\papereqref}[2]{(\ref{#1-#2})}

% CPSV16 source (1606.00608 / MPDO-22-12-17-2.tex) external document and shortcuts
\importpaperaux{cpsv16-}{1606.00608/MPDO-22-12-17-2-tnlean}
\newcommand{\cpsvref}[1]{\paperref{cpsv16}{#1}}
\newcommand{\cpsveqref}[1]{\papereqref{cpsv16}{#1}}

% Verdict marker: \gapnote{<kind>}{<status>}. Kind is the policy
% classification (clarification, local-correction, scope-restriction,
% unfaithful, false-source, open-gap); status is open, wip (elimination
% underway), resolved, or historical. Machine-read by the site index;
% prints a verdict line.
\newcommand{\gapnote}[2]{\par\noindent\textbf{Verdict:} #1 (#2).\par}


\title{A Phase Obstruction to Canonical-Form Renormalization-Flow Convergence}
\author{}
\date{2026-07-30}

\begin{document}
\maketitle
\gapnote{false-source}{resolved}

\section{The source assertion}

Around Theorem~3.1, Ref.~\cite{Cirac2016MPDO_arXiv} asserts at lines~382 and
417 that the renormalization flow starting from a tensor in canonical form
always converges. Appendix~B identifies one renormalization step with
\begin{align}
    \mathcal E
        &\longmapsto
    \mathcal E^2.
    \label{eq:cpsv16_canonical_form_renormalization_flow_phase_gap_renormalization_step}
\end{align}
This identification appears at lines~1205--1209; the canonical-form transfer
decomposition follows at lines~1211--1243. The conclusion at line~1244 is that
\begin{align}
    |\mu_{j,q}|
        &\leq
    1
    \label{eq:cpsv16_canonical_form_renormalization_flow_phase_gap_weight_normalization}
\end{align}
implies convergence of the transfer powers
\cite{Cirac2016MPDO_arXiv}.

This conclusion overlooks relative phases between repeated copies of the same
normal tensor. In the decomposition at lines~1220--1226, the corresponding
summands carry the factors
\begin{align}
    \mu_{j,q}\overline{\mu_{j,q'}}
    \mathbb E_{j,j}.
    \label{eq:cpsv16_canonical_form_renormalization_flow_phase_gap_repeated_copy_factor}
\end{align}
When both weights have modulus one,
Eq.~\ref{eq:cpsv16_canonical_form_renormalization_flow_phase_gap_weight_normalization}
does not force their relative phase to converge under repeated squaring.

\section{A phase-oscillation counterexample}

Define
\begin{align}
    \omega
        &=
    e^{2\pi i/3},
    \label{eq:cpsv16_canonical_form_renormalization_flow_phase_gap_third_root}
\end{align}
and consider the tensor with one physical component and bond dimension two,
\begin{align}
    A^1
        &=
    \begin{pmatrix}
        1 & 0\\
        0 & \omega
    \end{pmatrix}.
    \label{eq:cpsv16_canonical_form_renormalization_flow_phase_gap_counterexample_tensor}
\end{align}
This tensor satisfies the literal canonical-form definition at lines~238--246
of Ref.~\cite{Cirac2016MPDO_arXiv}: its two one-dimensional blocks are normal
tensors with normalized weights $1$ and $\omega$.

Let $e_{21}$ be the matrix unit with its only nonzero entry in row $2$ and
column $1$. The transfer map satisfies
\begin{align}
    \mathcal E_A(e_{21})
        &=
    A^1e_{21}(A^1)^*
    \label{eq:cpsv16_canonical_form_renormalization_flow_phase_gap_transfer_action_expansion}\\
        &=
    \omega e_{21}.
    \label{eq:cpsv16_canonical_form_renormalization_flow_phase_gap_transfer_eigenvalue}
\end{align}
Consequently,
\begin{align}
    \mathcal E_A^{\,2^n}(e_{21})
        &=
    \omega^{2^n}e_{21}.
    \label{eq:cpsv16_canonical_form_renormalization_flow_phase_gap_dyadic_transfer_power}
\end{align}
The residue of $2^n$ modulo $3$ alternates between $1$ and $2$. Hence the
sequence in
Eq.~\ref{eq:cpsv16_canonical_form_renormalization_flow_phase_gap_dyadic_transfer_power}
alternates between
\begin{align}
    \omega e_{21}
        \quad\text{and}\quad
    \omega^2e_{21}.
    \label{eq:cpsv16_canonical_form_renormalization_flow_phase_gap_alternating_values}
\end{align}
It does not converge. The unrestricted sequence $\mathcal E_A^N$ is likewise
three-periodic.

\section{Faithful boundary}

Theorem~3.1 of Ref.~\cite{Cirac2016MPDO_arXiv} remains valid: it characterizes
tensor classes that occur as limits and does not assert that every
canonical-form initial tensor has a limit. Primitive-transfer convergence for
each individual normal block also remains valid.

For a canonical form with repeated copies, convergence requires control of
their relative phases. A sufficient condition is that all unit-modulus weights
associated with one normal tensor have the same phase. More precisely, dyadic
convergence requires convergence of every sequence
\begin{align}
    \left(\mu_{j,q}\overline{\mu_{j,q'}}\right)^{2^n}
    \label{eq:cpsv16_canonical_form_renormalization_flow_phase_gap_relative_phase_sequence}
\end{align}
arising from two unit-modulus repeated-copy weights. Without such a condition,
the universal convergence assertion at lines~1211--1244 of
Ref.~\cite{Cirac2016MPDO_arXiv} is false rather than an unproved
formalization target.

\section{Conclusion}

The normalization in
Eq.~\ref{eq:cpsv16_canonical_form_renormalization_flow_phase_gap_weight_normalization}
controls the moduli of the canonical weights but not their relative phases.
The tensor in
Eq.~\ref{eq:cpsv16_canonical_form_renormalization_flow_phase_gap_counterexample_tensor}
satisfies the printed canonical-form assumptions and produces the
nonconvergent dyadic sequence in
Eq.~\ref{eq:cpsv16_canonical_form_renormalization_flow_phase_gap_dyadic_transfer_power}.
Thus the universal convergence claim in Appendix~B of
Ref.~\cite{Cirac2016MPDO_arXiv} is false as printed.

\bibliographystyle{alpha}
\bibliography{references}

\end{document}
